\documentclass[11pt,a4paper]{article}
\usepackage[utf8]{inputenc}
\usepackage[T1]{fontenc}
\usepackage[catalan]{babel}
\usepackage{lmodern}
\usepackage[sfdefault,lining]{FiraSans}
\usepackage[a4paper,top=2.3cm,bottom=2.2cm,left=1.8cm,right=1.8cm,headheight=1.3cm,headsep=0.4cm,footskip=1cm]{geometry}
\usepackage{amsmath,amssymb}
\usepackage{icomma}
\usepackage{graphicx}
\usepackage[table]{xcolor}
\usepackage{tikz}
\usepackage{fontawesome5}
\usepackage{ccicons}
\usepackage{fancyhdr}
\usepackage{lastpage}
\usepackage{tasks}
\usepackage{needspace}
\usepackage[most]{tcolorbox}
\usepackage{array}
\usepackage[hidelinks,pdftitle={Fitxa d'exercicis de la unitat 5: Sistemes d'equacions i inequacions. 4t d'ESO Matemàtiques A},pdfsubject={Matemàtiques A 4t ESO},pdfkeywords={sistemes d'equacions, mètode gràfic, substitució, igualació, reducció, parèntesis, denominadors, sistemes no lineals, problemes, inequacions de primer grau, inequacions de segon grau, sistemes d'inequacions}]{hyperref}

\definecolor{navy}{HTML}{1B2748}
\definecolor{accent}{HTML}{E8603A}
\definecolor{teal}{HTML}{1F9E90}
\definecolor{soft}{HTML}{EEF1F8}
\definecolor{muted}{HTML}{5D6685}

\setlength{\parindent}{0pt}
\setlength{\parskip}{3pt}
\settasks{label=\alph*), label-format=\bfseries\color{navy}, label-width=1.3em, item-indent=1.8em, column-sep=0.6em, after-item-skip=5pt, before-skip=3pt, after-skip=4pt}

% ---------- capçalera i peu ----------
\newcommand{\llicencia}{\href{https://creativecommons.org/licenses/by-nc-sa/4.0/deed.ca}{\ccbyncsa}\ \ Llicència Creative Commons Reconeixement-NoComercial-CompartirIgual 4.0 Internacional (CC BY-NC-SA 4.0)}
\pagestyle{fancy}
\fancyhf{}
\fancyhead[L]{\raisebox{-0.35\height}{\includegraphics[height=0.95cm]{logo-gva.png}}}
\fancyhead[C]{\small\color{muted}Fitxa d'exercicis · Unitat 5: Sistemes d'equacions i inequacions · 4t ESO A}
\fancyhead[R]{\raisebox{-0.35\height}{\includegraphics[height=0.95cm]{logo-vermellar.png}}}
\fancyfoot[L]{\scriptsize\color{muted}\llicencia}
\fancyfoot[R]{\small\color{muted}\thepage}
\renewcommand{\headrulewidth}{0.4pt}
\renewcommand{\headrule}{\color{muted!40}\hrule height \headrulewidth}
\fancypagestyle{primera}{\fancyhf{}\renewcommand{\headrulewidth}{0pt}
  \fancyfoot[L]{\scriptsize\color{muted}\llicencia}
  \fancyfoot[R]{\small\color{muted}\thepage}}

% ---------- nivells i exercicis ----------
\newcommand{\estrelles}[1]{{\color{accent}\ifcase#1\or\faStar\faStar[regular]\faStar[regular]\or\faStar\faStar\faStar[regular]\or\faStar\faStar\faStar\fi}}
\newcounter{exer}
\newcommand{\ex}[1]{\par\needspace{4\baselineskip}\medskip\refstepcounter{exer}\noindent{\color{navy}\bfseries\theexer.}\ \ {\small\estrelles{#1}}\ \ }

\newcommand{\seccio}[2]{\par\needspace{16\baselineskip}\bigskip
  \begin{tcolorbox}[colback=navy,colframe=navy,coltext=white,arc=2mm,boxrule=0pt,left=3mm,right=3mm,top=1.2mm,bottom=1.2mm]
  {\large\bfseries #1\quad #2}
  \end{tcolorbox}\vspace{-2pt}}

\newcommand{\exf}[1]{\par\medskip\refstepcounter{exer}\noindent{\color{navy}\bfseries\theexer.}\ \ {\small\estrelles{#1}}\ \ }
\newcolumntype{C}[1]{>{\centering\arraybackslash}p{#1}}
% quadre de recordatori (només a la fitxa d'exercicis)
\newcommand{\recorda}[1]{\par\smallskip\begin{tcolorbox}[colback=accent!7,colframe=accent!55,arc=2mm,boxrule=0.5pt,left=2.5mm,right=2.5mm,top=1mm,bottom=1mm]\small{\color{accent}\faLightbulb[regular]}\ \textbf{Recorda:}\ #1\end{tcolorbox}}
\usetikzlibrary{calc}
\newcommand{\res}[1]{#1}\newcommand{\rest}[1]{#1}

\def\fmneed{11}% que el títol d'una secció no quede sol al final del full
% =====================================================================
% Capçalera FentMates de les fitxes: un color per a cada curs.
% Els colors són els de la pàgina inicial (LEVELS a assets/js/materials.js);
% el logo ve de _eines/logo/pdf/ (python3 _eines/logo/genera.py).
%
% Ús, al final del preàmbul de la unitat:
%   % =====================================================================
% Capçalera FentMates de les fitxes: un color per a cada curs.
% Els colors són els de la pàgina inicial (LEVELS a assets/js/materials.js);
% el logo ve de _eines/logo/pdf/ (python3 _eines/logo/genera.py).
%
% Ús, al final del preàmbul de la unitat:
%   % =====================================================================
% Capçalera FentMates de les fitxes: un color per a cada curs.
% Els colors són els de la pàgina inicial (LEVELS a assets/js/materials.js);
% el logo ve de _eines/logo/pdf/ (python3 _eines/logo/genera.py).
%
% Ús, al final del preàmbul de la unitat:
%   \input{../comu/fentmates.tex}
%   \fmcurs{1eso}
% i, al començament del document:
%   \capcalera{Matemàtiques}{1r ESO}{Fitxa d'exercicis · Unitat 2}{Divisibilitat}
% Opcional abans de \input: \def\fmneed{14} (espai mínim abans de cada secció).
% =====================================================================
\makeatletter
\def\fm@logodir{../../logo/pdf}
\providecommand{\fmneed}{10}
\def\fm@cursnom{}

% clau del curs, color i color del text damunt del color
\newcommand{\fm@curs}[3]{\@namedef{fm@color@#1}{#2}\@namedef{fm@text@#1}{#3}}
\fm@curs{1eso}{E6B84A}{navy}
\fm@curs{2eso}{3B82F6}{white}
\fm@curs{3eso}{FF7A50}{white}
\fm@curs{4eso-a}{A78BFA}{white}
\fm@curs{4eso-b}{6D28D9}{white}
\fm@curs{1batx-mi}{9B1C1C}{white}
\fm@curs{1batx-ccss}{EF6B6B}{white}
\fm@curs{2batx-mii}{1E6B47}{white}
\fm@curs{2batx-ccss}{2BB3A3}{white}

\newcommand{\fmcurs}[1]{%
  \@ifundefined{fm@color@#1}{\PackageError{fentmates}{Curs desconegut: #1}{}}{}%
  \edef\fm@tmp{\noexpand\definecolor{curs}{HTML}{\@nameuse{fm@color@#1}}%
               \noexpand\colorlet{curstext}{\@nameuse{fm@text@#1}}}\fm@tmp
  \colorlet{cursfosc}{curs!72!black}%
  \def\fm@logo{\fm@logodir/vertical-#1.pdf}%
  \def\fm@marca{\fm@logodir/marca-#1.pdf}}

% ---------- capçalera de la primera pàgina ----------
% #1 matèria (petit), #2 curs, #3 tipus de document i unitat, #4 títol
\newcommand{\capcalera}[4]{%
  \gdef\fm@cursnom{#2}%
  \par\noindent
  \begin{tcolorbox}[colback=curs,colframe=curs,boxrule=0pt,arc=2.2mm,width=3.5cm,height=2.15cm,
      valign=center,left=3mm,right=2mm,top=1mm,bottom=1mm,nobeforeafter,box align=center]
    \color{curstext}{\footnotesize #1}\par\vspace{1pt}{\fontsize{19}{21}\selectfont\bfseries #2\par}%
  \end{tcolorbox}\hspace{5mm}%
  \begin{minipage}[c]{\dimexpr\linewidth-3.5cm-5mm-2.9cm\relax}
    {\color{cursfosc}\small\bfseries\MakeUppercase{#3}}\par\vspace{2pt}
    {\color{navy}\fontsize{20}{23}\selectfont\bfseries #4\par}
  \end{minipage}\hfill
  \begin{minipage}[c]{2.6cm}\raggedleft\includegraphics[width=2.6cm]{\fm@logo}\end{minipage}\par
  \vspace{2pt}}

% ---------- barres de secció del color del curs ----------
\renewcommand{\seccio}[2]{\par\needspace{\fmneed\baselineskip}\bigskip
  \begin{tcolorbox}[colback=curs,colframe=curs,coltext=curstext,arc=2mm,boxrule=0pt,left=3mm,right=3mm,top=1.2mm,bottom=1.2mm]
  {\large\bfseries #1\quad #2}
  \end{tcolorbox}\vspace{-2pt}}

% ---------- estrelles de dificultat del color del curs (a Batxillerat no se'n posen) ----------
\@ifundefined{estrelles}{}{%
  \renewcommand{\estrelles}[1]{{\color{curs}\ifcase#1\or\faStar\faStar[regular]\faStar[regular]\or\faStar\faStar\faStar[regular]\or\faStar\faStar\faStar\fi}}}

% ---------- capçalera i peu de les pàgines ----------
% primera pàgina: logos institucionals al peu
\fancypagestyle{primera}{\fancyhf{}\renewcommand{\headrulewidth}{0pt}
  \fancyfoot[L]{\raisebox{-0.32\height}{\includegraphics[height=0.78cm]{logo-gva.png}}\hspace{5mm}%
    \raisebox{-0.32\height}{\includegraphics[height=0.78cm]{logo-vermellar.png}}}
  \fancyfoot[R]{\scriptsize\color{muted}\href{https://creativecommons.org/licenses/by-nc-sa/4.0/deed.ca}{\ccbyncsa}\ \ CC BY-NC-SA 4.0 · fentmates.org\qquad\small\thepage}}
% la resta: el símbol a l'esquerra i el curs a la dreta, amb el títol al centre com abans
\fancyhead[L]{\raisebox{-0.3\height}{\includegraphics[height=0.62cm]{\fm@marca}}}
\fancyhead[R]{\small\bfseries\color{cursfosc}\fm@cursnom}
\renewcommand{\headrule}{\color{curs!45}\hrule height \headrulewidth}
\makeatother

%   \fmcurs{1eso}
% i, al començament del document:
%   \capcalera{Matemàtiques}{1r ESO}{Fitxa d'exercicis · Unitat 2}{Divisibilitat}
% Opcional abans de \input: \def\fmneed{14} (espai mínim abans de cada secció).
% =====================================================================
\makeatletter
\def\fm@logodir{../../logo/pdf}
\providecommand{\fmneed}{10}
\def\fm@cursnom{}

% clau del curs, color i color del text damunt del color
\newcommand{\fm@curs}[3]{\@namedef{fm@color@#1}{#2}\@namedef{fm@text@#1}{#3}}
\fm@curs{1eso}{E6B84A}{navy}
\fm@curs{2eso}{3B82F6}{white}
\fm@curs{3eso}{FF7A50}{white}
\fm@curs{4eso-a}{A78BFA}{white}
\fm@curs{4eso-b}{6D28D9}{white}
\fm@curs{1batx-mi}{9B1C1C}{white}
\fm@curs{1batx-ccss}{EF6B6B}{white}
\fm@curs{2batx-mii}{1E6B47}{white}
\fm@curs{2batx-ccss}{2BB3A3}{white}

\newcommand{\fmcurs}[1]{%
  \@ifundefined{fm@color@#1}{\PackageError{fentmates}{Curs desconegut: #1}{}}{}%
  \edef\fm@tmp{\noexpand\definecolor{curs}{HTML}{\@nameuse{fm@color@#1}}%
               \noexpand\colorlet{curstext}{\@nameuse{fm@text@#1}}}\fm@tmp
  \colorlet{cursfosc}{curs!72!black}%
  \def\fm@logo{\fm@logodir/vertical-#1.pdf}%
  \def\fm@marca{\fm@logodir/marca-#1.pdf}}

% ---------- capçalera de la primera pàgina ----------
% #1 matèria (petit), #2 curs, #3 tipus de document i unitat, #4 títol
\newcommand{\capcalera}[4]{%
  \gdef\fm@cursnom{#2}%
  \par\noindent
  \begin{tcolorbox}[colback=curs,colframe=curs,boxrule=0pt,arc=2.2mm,width=3.5cm,height=2.15cm,
      valign=center,left=3mm,right=2mm,top=1mm,bottom=1mm,nobeforeafter,box align=center]
    \color{curstext}{\footnotesize #1}\par\vspace{1pt}{\fontsize{19}{21}\selectfont\bfseries #2\par}%
  \end{tcolorbox}\hspace{5mm}%
  \begin{minipage}[c]{\dimexpr\linewidth-3.5cm-5mm-2.9cm\relax}
    {\color{cursfosc}\small\bfseries\MakeUppercase{#3}}\par\vspace{2pt}
    {\color{navy}\fontsize{20}{23}\selectfont\bfseries #4\par}
  \end{minipage}\hfill
  \begin{minipage}[c]{2.6cm}\raggedleft\includegraphics[width=2.6cm]{\fm@logo}\end{minipage}\par
  \vspace{2pt}}

% ---------- barres de secció del color del curs ----------
\renewcommand{\seccio}[2]{\par\needspace{\fmneed\baselineskip}\bigskip
  \begin{tcolorbox}[colback=curs,colframe=curs,coltext=curstext,arc=2mm,boxrule=0pt,left=3mm,right=3mm,top=1.2mm,bottom=1.2mm]
  {\large\bfseries #1\quad #2}
  \end{tcolorbox}\vspace{-2pt}}

% ---------- estrelles de dificultat del color del curs (a Batxillerat no se'n posen) ----------
\@ifundefined{estrelles}{}{%
  \renewcommand{\estrelles}[1]{{\color{curs}\ifcase#1\or\faStar\faStar[regular]\faStar[regular]\or\faStar\faStar\faStar[regular]\or\faStar\faStar\faStar\fi}}}

% ---------- capçalera i peu de les pàgines ----------
% primera pàgina: logos institucionals al peu
\fancypagestyle{primera}{\fancyhf{}\renewcommand{\headrulewidth}{0pt}
  \fancyfoot[L]{\raisebox{-0.32\height}{\includegraphics[height=0.78cm]{logo-gva.png}}\hspace{5mm}%
    \raisebox{-0.32\height}{\includegraphics[height=0.78cm]{logo-vermellar.png}}}
  \fancyfoot[R]{\scriptsize\color{muted}\href{https://creativecommons.org/licenses/by-nc-sa/4.0/deed.ca}{\ccbyncsa}\ \ CC BY-NC-SA 4.0 · fentmates.org\qquad\small\thepage}}
% la resta: el símbol a l'esquerra i el curs a la dreta, amb el títol al centre com abans
\fancyhead[L]{\raisebox{-0.3\height}{\includegraphics[height=0.62cm]{\fm@marca}}}
\fancyhead[R]{\small\bfseries\color{cursfosc}\fm@cursnom}
\renewcommand{\headrule}{\color{curs!45}\hrule height \headrulewidth}
\makeatother

%   \fmcurs{1eso}
% i, al començament del document:
%   \capcalera{Matemàtiques}{1r ESO}{Fitxa d'exercicis · Unitat 2}{Divisibilitat}
% Opcional abans de \input: \def\fmneed{14} (espai mínim abans de cada secció).
% =====================================================================
\makeatletter
\def\fm@logodir{../../logo/pdf}
\providecommand{\fmneed}{10}
\def\fm@cursnom{}

% clau del curs, color i color del text damunt del color
\newcommand{\fm@curs}[3]{\@namedef{fm@color@#1}{#2}\@namedef{fm@text@#1}{#3}}
\fm@curs{1eso}{E6B84A}{navy}
\fm@curs{2eso}{3B82F6}{white}
\fm@curs{3eso}{FF7A50}{white}
\fm@curs{4eso-a}{A78BFA}{white}
\fm@curs{4eso-b}{6D28D9}{white}
\fm@curs{1batx-mi}{9B1C1C}{white}
\fm@curs{1batx-ccss}{EF6B6B}{white}
\fm@curs{2batx-mii}{1E6B47}{white}
\fm@curs{2batx-ccss}{2BB3A3}{white}

\newcommand{\fmcurs}[1]{%
  \@ifundefined{fm@color@#1}{\PackageError{fentmates}{Curs desconegut: #1}{}}{}%
  \edef\fm@tmp{\noexpand\definecolor{curs}{HTML}{\@nameuse{fm@color@#1}}%
               \noexpand\colorlet{curstext}{\@nameuse{fm@text@#1}}}\fm@tmp
  \colorlet{cursfosc}{curs!72!black}%
  \def\fm@logo{\fm@logodir/vertical-#1.pdf}%
  \def\fm@marca{\fm@logodir/marca-#1.pdf}}

% ---------- capçalera de la primera pàgina ----------
% #1 matèria (petit), #2 curs, #3 tipus de document i unitat, #4 títol
\newcommand{\capcalera}[4]{%
  \gdef\fm@cursnom{#2}%
  \par\noindent
  \begin{tcolorbox}[colback=curs,colframe=curs,boxrule=0pt,arc=2.2mm,width=3.5cm,height=2.15cm,
      valign=center,left=3mm,right=2mm,top=1mm,bottom=1mm,nobeforeafter,box align=center]
    \color{curstext}{\footnotesize #1}\par\vspace{1pt}{\fontsize{19}{21}\selectfont\bfseries #2\par}%
  \end{tcolorbox}\hspace{5mm}%
  \begin{minipage}[c]{\dimexpr\linewidth-3.5cm-5mm-2.9cm\relax}
    {\color{cursfosc}\small\bfseries\MakeUppercase{#3}}\par\vspace{2pt}
    {\color{navy}\fontsize{20}{23}\selectfont\bfseries #4\par}
  \end{minipage}\hfill
  \begin{minipage}[c]{2.6cm}\raggedleft\includegraphics[width=2.6cm]{\fm@logo}\end{minipage}\par
  \vspace{2pt}}

% ---------- barres de secció del color del curs ----------
\renewcommand{\seccio}[2]{\par\needspace{\fmneed\baselineskip}\bigskip
  \begin{tcolorbox}[colback=curs,colframe=curs,coltext=curstext,arc=2mm,boxrule=0pt,left=3mm,right=3mm,top=1.2mm,bottom=1.2mm]
  {\large\bfseries #1\quad #2}
  \end{tcolorbox}\vspace{-2pt}}

% ---------- estrelles de dificultat del color del curs (a Batxillerat no se'n posen) ----------
\@ifundefined{estrelles}{}{%
  \renewcommand{\estrelles}[1]{{\color{curs}\ifcase#1\or\faStar\faStar[regular]\faStar[regular]\or\faStar\faStar\faStar[regular]\or\faStar\faStar\faStar\fi}}}

% ---------- capçalera i peu de les pàgines ----------
% primera pàgina: logos institucionals al peu
\fancypagestyle{primera}{\fancyhf{}\renewcommand{\headrulewidth}{0pt}
  \fancyfoot[L]{\raisebox{-0.32\height}{\includegraphics[height=0.78cm]{logo-gva.png}}\hspace{5mm}%
    \raisebox{-0.32\height}{\includegraphics[height=0.78cm]{logo-vermellar.png}}}
  \fancyfoot[R]{\scriptsize\color{muted}\href{https://creativecommons.org/licenses/by-nc-sa/4.0/deed.ca}{\ccbyncsa}\ \ CC BY-NC-SA 4.0 · fentmates.org\qquad\small\thepage}}
% la resta: el símbol a l'esquerra i el curs a la dreta, amb el títol al centre com abans
\fancyhead[L]{\raisebox{-0.3\height}{\includegraphics[height=0.62cm]{\fm@marca}}}
\fancyhead[R]{\small\bfseries\color{cursfosc}\fm@cursnom}
\renewcommand{\headrule}{\color{curs!45}\hrule height \headrulewidth}
\makeatother

\fmcurs{4eso-a}

\begin{document}
\thispagestyle{primera}

\capcalera{Matemàtiques A}{4t ESO}{Fitxa d'exercicis · Unitat 5}{\raggedright Sistemes d'equacions i inequacions}
\vspace{10pt}
Nom i cognoms: \rule{8.4cm}{0.4pt}\hfill Grup: \rule{1.6cm}{0.4pt}\hfill Data: \rule{2.2cm}{0.4pt}
\vspace{8pt}

\begin{tcolorbox}[colback=soft,colframe=soft,arc=2mm,left=3mm,right=3mm,top=1.5mm,bottom=1.5mm]
\small Exercicis graduats per nivells de dificultat:\quad \estrelles{1}\ bàsic \qquad \estrelles{2}\ mitjà \qquad \estrelles{3}\ avançat.\\[2pt] Resol-los al quadern mostrant tots els passos. Comprova les solucions dels sistemes en les dues equacions i, en els problemes, respon amb una frase.
\end{tcolorbox}

% =====================================================
\seccio{1}{Sistemes lineals. Mètode gràfic}

\exf{1} Comprova si $(2,-1)$ i $(1,\,1)$ són solució del sistema $\begin{cases}3x+2y=4\\[2pt] x-y=3\end{cases}$.

\ex{1} Troba quatre solucions de l'equació $2x+y=6$ amb una taula de valors i representa-les. Què formen?

\ex{2} Resol gràficament els sistemes:\begin{tasks}[after-item-skip=9pt](3)
  \task $\begin{cases}x+y=5\\[2pt] x-y=1\end{cases}$
  \task $\begin{cases}2x-y=1\\[2pt] x+y=5\end{cases}$
  \task $\begin{cases}x-2y=-4\\[2pt] x+y=2\end{cases}$
\end{tasks}\recorda{Fes una taula de valors amb dos punts de cada recta. La solució és el punt de tall.}

\ex{2} Classifica els sistemes sense resoldre'ls:\begin{tasks}[after-item-skip=9pt](3)
  \task $\begin{cases}2x+3y=5\\[2pt] 4x+6y=10\end{cases}$
  \task $\begin{cases}x-y=2\\[2pt] 3x-3y=1\end{cases}$
  \task $\begin{cases}x+2y=3\\[2pt] 2x-y=1\end{cases}$
\end{tasks}

\ex{3} Troba el valor de $k$ perquè el sistema $\begin{cases}kx+2y=4\\[2pt] 3x+y=5\end{cases}$ siga incompatible.


% =====================================================
\seccio{2}{Mètodes de substitució, igualació i reducció}

\exf{1} Resol pel mètode de substitució:\begin{tasks}[after-item-skip=9pt](4)
  \task $\begin{cases}y=2x-3\\[2pt] 3x+y=7\end{cases}$
  \task $\begin{cases}x+3y=10\\[2pt] 2x-y=6\end{cases}$
  \task $\begin{cases}x=3y-1\\[2pt] 2x+y=12\end{cases}$
  \task $\begin{cases}y=5-x\\[2pt] 3x-2y=5\end{cases}$
\end{tasks}

\ex{1} Resol pel mètode d'igualació:\begin{tasks}[after-item-skip=9pt](4)
  \task $\begin{cases}y=x+4\\[2pt] y=3x\end{cases}$
  \task $\begin{cases}x=2y-1\\[2pt] x=8-y\end{cases}$
  \task $\begin{cases}y=2x-1\\[2pt] y=-x+5\end{cases}$
  \task $\begin{cases}x=4-y\\[2pt] x=2y+1\end{cases}$
\end{tasks}

\ex{2} Resol pel mètode de reducció:\begin{tasks}[after-item-skip=9pt](4)
  \task $\begin{cases}2x+3y=12\\[2pt] 2x-y=4\end{cases}$
  \task $\begin{cases}3x+2y=7\\[2pt] 5x-3y=18\end{cases}$
  \task $\begin{cases}4x+3y=1\\[2pt] 3x-2y=5\end{cases}$
  \task $\begin{cases}5x+2y=1\\[2pt] 3x-4y=11\end{cases}$
\end{tasks}\recorda{Multiplica les equacions perquè una incògnita tinga el mateix coeficient (o l'oposat) i resta-les (o suma-les).}

\ex{2} Resol el sistema $\begin{cases}x+y=9\\[2pt] x-y=3\end{cases}$ pels tres mètodes algebraics i comprova que obtens la mateixa solució.

\ex{2} Resol i classifica:\begin{tasks}[after-item-skip=9pt](2)
  \task $\begin{cases}2x-4y=6\\[2pt] -x+2y=-3\end{cases}$
  \task $\begin{cases}x+y=3\\[2pt] 2x+2y=7\end{cases}$
\end{tasks}


% =====================================================
\seccio{3}{Sistemes amb parèntesis i denominadors}

\exf{2} Resol (primer trau els parèntesis):\begin{tasks}[after-item-skip=9pt](3)
  \task $\begin{cases}3(x-1)+y=6\\[2pt] x-2(y+1)=-6\end{cases}$
  \task $\begin{cases}2(x+1)-y=7\\[2pt] x+3(y-2)=0\end{cases}$
  \task $\begin{cases}5(x-y)+2y=13\\[2pt] 2(x+y)-3x=-4\end{cases}$
\end{tasks}

\ex{2} Resol (primer elimina els denominadors):\begin{tasks}[after-item-skip=16pt](3)
  \task $\begin{cases}\dfrac{x}{3}+\dfrac{y}{2}=4\\[9pt] \dfrac{x}{2}-\dfrac{y}{4}=2\end{cases}$
  \task $\begin{cases}\dfrac{x}{2}+\dfrac{y}{3}=3\\[9pt] x-y=1\end{cases}$
  \task $\begin{cases}\dfrac{x}{2}+y=5\\[9pt] \dfrac{x+2}{3}-\dfrac{y-1}{2}=1\end{cases}$
\end{tasks}\recorda{Multiplica cada equació pel m.c.m. dels seus denominadors i reduïx-la a la forma $ax+by=c$.}

\ex{3} Resol:\begin{tasks}[after-item-skip=16pt](2)
  \task $\begin{cases}\dfrac{x+1}{2}+\dfrac{y}{3}=2\\[9pt] 2(x-y)=x-3\end{cases}$
  \task $\begin{cases}\dfrac{x-1}{2}+\dfrac{y+1}{3}=2\\[9pt] 2(x+y)-(x-3)=10\end{cases}$
  \task $\begin{cases}\dfrac{x-y}{2}+y=4\\[9pt] \dfrac{x+y}{3}-\dfrac{x-y}{6}=1\end{cases}$
  \task $\begin{cases}\dfrac{x}{5}-\dfrac{y}{2}=0\\[9pt] \dfrac{x+1}{3}+\dfrac{y-2}{4}=2\end{cases}$
\end{tasks}

\ex{3} Resol:\begin{tasks}[after-item-skip=16pt](3)
  \task $\begin{cases}\dfrac{x+2}{3}=\dfrac{y+1}{2}\\[9pt] 2x-y=5\end{cases}$
  \task $\begin{cases}\dfrac{x+y}{2}+\dfrac{x-y}{4}=\dfrac{5}{2}\\[9pt] 2(x+1)-3(y-1)=8\end{cases}$
  \task $\begin{cases}\dfrac{x-1}{3}+\dfrac{y+3}{2}=\dfrac{4}{3}\\[9pt] 3(x+y)-2(x-y)=-3\end{cases}$
\end{tasks}


% =====================================================
\seccio{4}{Sistemes no lineals}

\exf{2} Resol pel mètode de substitució:\begin{tasks}[after-item-skip=9pt](3)
  \task $\begin{cases}x+y=5\\[2pt] x^{2}+y^{2}=13\end{cases}$
  \task $\begin{cases}x-y=2\\[2pt] xy=15\end{cases}$
  \task $\begin{cases}y=x^{2}-4\\[2pt] y=2x-1\end{cases}$
\end{tasks}

\ex{2} Resol pel mètode de reducció:\begin{tasks}[after-item-skip=9pt](2)
  \task $\begin{cases}x^{2}+y^{2}=25\\[2pt] x^{2}-y^{2}=7\end{cases}$
  \task $\begin{cases}x^{2}+2y^{2}=6\\[2pt] 3x^{2}-y^{2}=11\end{cases}$
\end{tasks}

\ex{2} Resol:\begin{tasks}[after-item-skip=9pt](4)
  \task $\begin{cases}x^{2}+y^{2}=10\\[2pt] x+y=4\end{cases}$
  \task $\begin{cases}y=x^{2}-2x\\[2pt] y=x+4\end{cases}$
  \task $\begin{cases}x-y=1\\[2pt] xy=6\end{cases}$
  \task $\begin{cases}x+y=6\\[2pt] x^{2}-y=6\end{cases}$
\end{tasks}

\ex{2} Sense dibuixar, esbrina si la recta $y=x+3$ talla la circumferència $x^2+y^2=4$.


% =====================================================
\seccio{5}{Problemes amb sistemes lineals}

\exf{1} En una classe hi ha 28 alumnes, i el nombre de xiques supera en 4 el de xics. Quantes xiques i quants xics hi ha?

\ex{1} En un aparcament hi ha 30 vehicles entre cotxes i motos, i en total 100 rodes. Quants n'hi ha de cada tipus?

\ex{1} Tinc 25 monedes d'1 € i de 2 €, que fan 38 €. Quantes en tinc de cada classe?

\ex{2} Tres entrades d'adult i dues d'infant costen 33 €, i dues d'adult i tres d'infant, 27 €. Quant costa cada entrada?

\ex{2} Un pare té 30 anys més que el seu fill. D'ací a 10 anys, l'edat del pare serà el doble de la del fill. Quants anys té cadascun?

\ex{2} Les xifres d'un nombre de dues xifres sumen 11. Si les intercanviem, el nombre disminueix 27. Quin és el nombre?

\ex{3} Volem fer 20 kg d'una mescla de cafè a 9,40 €/kg amb cafè de 8 €/kg i cafè de 12 €/kg. Quants quilos n'hem de posar de cada classe?


% =====================================================
\seccio{6}{Problemes amb sistemes no lineals}

\exf{2} Un rectangle té 32 cm de perímetre i 60 cm$^2$ d'àrea. Quant fan els costats?

\ex{2} El perímetre d'un rectangle és de 28 m i la diagonal fa 10 m. Quant fan els costats?

\ex{2} Dos nombres es diferencien en 3 i la suma dels seus quadrats és 117. Quins són?

\ex{3} La hipotenusa d'un triangle rectangle fa 13 cm i l'àrea, 30 cm$^2$. Quant fan els catets?


% =====================================================
\seccio{7}{Inequacions de primer grau}

\exf{1} Resol i representa la solució:\begin{tasks}(4)
  \task $3x-7<5$
  \task $2x+5\ge x+3$
  \task $4-3x\ge 10$
  \task $5x+3>2x-6$
\end{tasks}\recorda{Si multipliques o divideixes per un nombre negatiu, la desigualtat canvia de sentit.}

\ex{2} Resol (primer trau els parèntesis):\begin{tasks}(2)
  \task $3(x-2)+1<x+5$
  \task $2-(x+4)\ge 3x+6$
  \task $2(x-3)-(x+1)>3x-5$
  \task $4(x+1)-3(2x-1)\le x+1$
\end{tasks}

\ex{2} Resol (primer elimina els denominadors):\begin{tasks}[after-item-skip=12pt](2)
  \task $\dfrac{x+1}{2}-\dfrac{x-3}{3}\le 1$
  \task $\dfrac{x}{4}-\dfrac{2x-1}{6}>\dfrac{x}{3}+1$
  \task $\dfrac{x-1}{2}-\dfrac{x+2}{3}\ge 1$
  \task $\dfrac{x+3}{4}-\dfrac{x-1}{6}\ge 1$
\end{tasks}

\ex{2} Resol i explica què passa:\begin{tasks}(2)
  \task $2(x+3)>2x+1$
  \task $3(x-1)<3x-5$
\end{tasks}

\ex{2} En els dos primers exàmens he tret un 4 i un 6,5. Quina nota he de traure com a mínim en el tercer perquè la mitjana siga almenys un 5?

\ex{3} La companyia A cobra 15 € al mes més 0,10 € per minut, i la B, 0,25 € per minut. A partir de quants minuts al mes és més barata la companyia A?


% =====================================================
\seccio{8}{Inequacions de segon grau}

\exf{2} Resol:\begin{tasks}(4)
  \task $x^{2}-5x+6<0$
  \task $x^{2}-9\ge 0$
  \task $x^{2}+3x\le 0$
  \task $x^{2}-x-6>0$
\end{tasks}\recorda{Calcula les arrels i estudia el signe en cada interval. Si $a>0$, la paràbola és negativa entre les arrels i positiva fora.}

\ex{2} Resol:\begin{tasks}(4)
  \task $-x^{2}+2x+3>0$
  \task $-x^{2}+4x-3\ge 0$
  \task $2x^{2}+6x>0$
  \task $3x^{2}-12\le 0$
\end{tasks}

\ex{3} Resol (primer trau els parèntesis):\begin{tasks}(2)
  \task $(x-1)^{2}\ge x+5$
  \task $x(x+2)<3$
  \task $(x+2)(x-2)>3x$
  \task $(x+1)^{2}<2(x+5)$
\end{tasks}

\ex{3} Resol (compte: el polinomi té una arrel doble o no en té):\begin{tasks}(4)
  \task $x^{2}-6x+9\le 0$
  \task $x^{2}+2x+1>0$
  \task $x^{2}+4>0$
  \task $x^{2}+2x+5<0$
\end{tasks}

\ex{3} Els costats d'un rectangle fan $x$ i $x+2$ cm. Per a quins valors de $x$ l'àrea és menor que 15 cm$^2$?


% =====================================================
\seccio{9}{Sistemes d'inequacions}

\exf{1} Resol els sistemes i representa la solució:\begin{tasks}[after-item-skip=9pt](4)
  \task $\begin{cases}x-3<0\\[2pt] x+1\ge 0\end{cases}$
  \task $\begin{cases}2x+1>5\\[2pt] 3x-2\le 10\end{cases}$
  \task $\begin{cases}x+4<1\\[2pt] 2x>-2\end{cases}$
  \task $\begin{cases}5-x\le 2\\[2pt] 4x+1\ge 9\end{cases}$
\end{tasks}\recorda{Resol cada inequació i fes la intersecció de les solucions: dibuixa-les una damunt de l'altra.}

\ex{2} Resol:\begin{tasks}[after-item-skip=16pt](3)
  \task $\begin{cases}3(x-1)\ge x+1\\[9pt] \dfrac{x+2}{3}<2\end{cases}$
  \task $\begin{cases}2x-5<x+1\\[9pt] \dfrac{x}{2}+1>0\end{cases}$
  \task $\begin{cases}2(x-1)\ge x-3\\[9pt] \dfrac{x+4}{2}>x\end{cases}$
\end{tasks}

\ex{3} Resol (alguna inequació és de segon grau):\begin{tasks}[after-item-skip=9pt](2)
  \task $\begin{cases}x^{2}-9\le 0\\[2pt] 2x-1>0\end{cases}$
  \task $\begin{cases}x^{2}-4x>0\\[2pt] x+2\ge 0\end{cases}$
\end{tasks}

\ex{2} Un taxi cobra 3 € per la baixada de bandera i 1,20 € per quilòmetre. Si volem gastar com a mínim 9 € i com a màxim 15 €, quants quilòmetres podem fer?

\ex{3} Troba els nombres naturals que compleixen alhora que el seu doble menys 3 és major que 7 i que el seu triple més 1 és menor que 25.

\vfill
\begin{center}\color{muted}\small\faCheckCircle[regular]\ \ Ànim! Comprova els teus resultats amb les eines interactives dels apunts de la unitat.\end{center}

\end{document}
